\section{Conclusões e trabalho futuro}
\label{sec:conclusoes}

Este trabalho apresentou uma revisão sistemática da literatura sobre segurança em
comunicações IoT baseadas em MQTT, com \nTriados{} registos triados e \nIncluidos{} estudos primários (2015--2026)
e \nSecundarios{} revisões relacionadas analisados em texto integral. As principais conclusões são:

\begin{itemize}[leftmargin=*]
  \item O TLS, por terminar no \emph{broker}, não protege as mensagens de um
    \emph{broker} comprometido nem de clientes com acesso legítimo.
  \item A proteção leve fim-a-fim é a categoria mais numerosa, e a
    confidencialidade, a autenticação e a integridade são as propriedades mais
    tratadas.
  \item A validade temporal, a deteção de duplicados, a ligação ao tópico e o não
    repúdio aparecem em poucos estudos e nunca combinados; a restrição
    chave--tópico no recetor e a política contra a redução do nível de proteção não
    aparecem nos estudos analisados; e a maioria das soluções confia no
    \emph{broker}.
  \item As avaliações medem o custo dos mecanismos, mas não o relacionam com os
    ataques impedidos.
\end{itemize}

Com base nestas lacunas, propôs-se, não um mecanismo novo, mas a integração destes
elementos numa arquitetura em camadas independentes e
comutáveis, com defesa no transporte, um envelope fim-a-fim (AES-256-GCM ou ChaCha20-Poly1305, GMAC ou
Ed25519, com tópico e cabeçalho nos dados associados, validade temporal,
anti-repetição, política de modos e restrição chave--tópico) e limitação de
tráfego antes da criptografia.

Como trabalho futuro, fica sobretudo a segunda parte do projeto
(Secção~\ref{sec:plano}): medir num Raspberry Pi o custo e os ataques impedidos por
cada configuração e escolher a melhor pela fronteira de Pareto. Seguem-se a
análise em texto integral dos restantes estudos elegíveis, a verificação formal do
envelope, à semelhança de~\cite{javadi2026cmap}, um acordo de
chaves por dispositivo e a integração com deteção de
intrusões~\cite{vaccari2020mqttset, alotaibi2024secure}.
